\bookchapter{الفصل 27: يستخدمه الآباء بسهولة — دليل شامل لهاتف مناسب لكبار السن}

\chapterintro{بصراحة، عند كتابة هذا الفصل، فكرت في والديّ. والدتي لا تجرؤ على النقر على “تحديث النظام” في كل مرة، خوفاً من إتلاف الهاتف. خط WeChat لأبي أصغر من النمل، لكنه لا يعرف كيف يكبره، وظل يضيق عينيه لسنتين. هذا ليس “ضعفاً” منهما، بل لأن الهاتف يفترض أن الجميع في العشرينات. هذا الفصل، بأوضح كلمة، يحل مشاكل كبار السن في استخدام الهاتف واحدة تلو الأخرى.}

\booksection{صفر: ضبط العقلية أولاً: ليست “كبار السن لا يستطيعون” بل “الهاتف لم يُصمم جيداً”}

قبل أي إعداد، يجب توضيح فكرة:\textbf{واجهات الهواتف الذكية الحالية لم تأخذ المستخدمين كبار السن بعين الاعتبار منذ التصميم}. النوافذ المنبثقة صغيرة لا تُرى بوضوح، والإيماءات تتطلب لمساً دقيقاً، والأيقونات متشابهة — هذه ليست مشاكل المستخدم بل مشاكل المنتج. لذا عند مساعدة كبار السن، أول شيء ليس التعليم بل مساعدتهم\textbf{ضبط البيئة}إذا ضُبطت البيئة، تختفي أغلب المشاكل.

ثلاثة قواعد ذهبية لمساعدة كبار السن: (1)\textbf{لا تتعجل}— أنت تستخدم الهاتف الذكي منذ عشر سنوات، هم ربما منذ ثلاثة أشهر. (2)\textbf{استخدم التشبيهات}— “Wi-Fi مثل أنبوب الماء في المنزل، افتح الصنبور ليأتي الماء”، هذا التشبيه أسهل فهماً بآلاف المرات من المصطلحات التقنية. (3)\textbf{اكتب}— قد ينسون بعد تعليم شفهي، خذ ورقة واكتب الخطوات الرئيسية والصقها على ظهر غطاء الهاتف، أنفع من أي شيء.

\booksection{واحد: الخطوة الأولى: تفعيل “الوضع البسيط” — الهاتف يصبح أنظف}

الهواتف الحالية تحتوي عموماً على “وضع بسيط” مدمج (يسمى أحياناً “وضع الرعاية” “وضع كبار السن” “وضع المسنين”)، بعد تفعيله، يتغير الهاتف بالكامل:

\begin{itemize}
  \item \textbf{الخطوط والأيقونات تصبح كبيرة جداً}— لم تعد بحاجة لنظارة قراءة.
  \item \textbf{أيقونات سطح المكتب مُبسّطة كثيراً}— الاحتفاظ فقط بالهاتف والرسائل والكاميرا والألبوم وWeChat وغيرها من التطبيقات الأكثر استخداماً، وإخفاء الباقي.
  \item \textbf{العمليات تصبح أبسط}— الكثير من الوظائف الاستعراضية (الشاشة السلبية والمساعد الذكي والنوافذ العائمة) تُغلق افتراضياً، والواجهة أساسها النظافة.
  \item \textbf{الصوت يصبح أعلى}— نغمة المكالمات الواردة وصوت الإشعارات وصوت المكالمات يُضبط تلقائياً على الأعلى.
\end{itemize}

\booksubsection{طرق تفعيل الوضع البسيط لكل علامة (اتبعها مباشرة):}

\begin{center}
\begin{tabularx}{\linewidth}{|X|X|X|}
\hline
\textbf{العلامة} & \textbf{اسم الوضع} & \textbf{طريقة التفعيل} \tabularnewline
\hline
iPhone & الوصول المساعد (لا وضع بسيط منفصل، لكن يمكن التكبير عنصراً بعنصر) & الإعدادات > العرض والسطوع > حجم النص، اسحب للأكبر. ثم "الوصول المساعد > العرض وحجم النص" فعّل "نص عريض" "أشكال الأزرار" \tabularnewline
Huawei/Honor & الوضع البسيط & الإعدادات > النظام والتحديث > الوضع البسيط > تفعيل \tabularnewline
Xiaomi/Redmi & الوضع المبسط & الإعدادات > إعدادات أكثر > الوضع المبسط > تفعيل \tabularnewline
OPPO/OnePlus & الوضع البسيط & الإعدادات > الميزات المميزة > الوضع البسيط \tabularnewline
vivo/iQOO & سطح المكتب المبسط & الإعدادات > سطح المكتب وشاشة القفل والخلفية > سطح المكتب المبسط \tabularnewline
Samsung & الوضع البسيط & الإعدادات > العرض > الوضع البسيط \tabularnewline
\hline
\end{tabularx}
\end{center}

إذا كانت علامتك أخرى، ابحث بكلمة “بسيط” أو “مسنين” في مربع البحث أعلى “الإعدادات”، وغالباً ستجده. إذا لم تجده، اتبع خطة “التكبير اليدوي” أدناه عنصراً عنصراً، النتيجة نفسها.

\booksection{اثنان: إذا لم تفعّل الوضع البسيط، يمكن التكبير يدوياً — “خط كبير + أيقونات كبيرة” بخطوات خمس}

بعض كبار السن لا يريدون الوضع البسيط، يرونه “عجوزاً”. يمكن التكبير اليدوي للخمسة التالية، نفس النتيجة:

\begin{enumerate}
  \item \textbf{تكبير الخط لأكبر حجم}: الإعدادات > العرض > حجم الخط، اسحب لليمين. واضبط "حجم العرض" (أو "تكبير العرض") أيضاً للأكبر، هكذا تُكبر جميع عناصر الواجهة.
  \item \textbf{تكبير الأيقونات}: على الشاشة الرئيسية قرّب إصبعين (مثل تصغير الصورة) للدخول لإعدادات سطح المكتب > تخطيط سطح المكتب (أو حجم الأيقونات)، اختر التخطيط بأقل خانات (مثل 3×5)، فتكبر الأيقونات.
  \item \textbf{تفعيل "نص عالي التباين"}: يجعل حواف النص أكثر حدة، يمكن إيجاده في "الوصول المساعد".
  \item \textbf{إغلاق الإشعارات غير الضرورية}: الإعدادات > الإشعارات، أغلق إشعارات جميع التطبيقات باستثناء الهاتف والرسائل وWeChat. هكذا لن يصدر الهاتف صوتاً كل لحظة ولن يزعج.
  \item \textbf{وضع جهات الاتصال المعتادة على سطح المكتب}: اضغط مطولاً على مكان فارغ في سطح المكتب > إضافة أداة > جهات الاتصال، ضع الاتصال السريع للزوج/الزوجة والأبناء والطبيب وجهة الطوارئ. نقرة واحدة للاتصال، دون فتح دليل العناوين.
\end{enumerate}

\begin{tcbpractice}{جرّب: ساعد كبار السن في "تزيين" الهاتف}
\begin{enumerate}
  \item افتح الوضع البسيط (اتبع جدول العلامات أعلاه).
  \item فعّل عرض نسبة البطارية (الإعدادات > البطارية > نسبة البطارية)، ليعرفوا البطارية المتبقية في أي وقت.
  \item اعرض شريط الإشعارات، أزل أو أخفِ مفاتيح Bluetooth وNFC ووضع الطيران التي لا يحتاجها كبار السن.
  \item ضع أربعة اتصالات سريعة على سطح المكتب: الزوج/الزوجة والأبناء وطبيب الأسرة (إن وجد) والمجتمع/الإدارة.
\end{enumerate}

\end{tcbpractice}

\booksection{ثلاثة: مكافحة الاحتيال — هذا القسم اقرأه مع كبار السن}

كبار السن أكثر عرضة للاحتيال الهاتفي. ليس لأن حذرهم منخفض، بل لأن أساليب المحتالين مصممة بعناية وتستهدف نقاط ضعف كبار السن النفسية. التالي أكثر أنواع الاحتيال شيوعاً وطرق التعامل، يُنصح بقراءتها لكبار السن أو طباعتها ولصقها بجانب الهاتف.

\booksubsection{ستة احتيالات شائعة (إذا رأيتها أغلق المكالمة فوراً):}

\begin{itemize}
  \item \textbf{(1) "حفيدك في مشكلة، أرسل المال بسرعة"}— يقول المحتال "حفيدك ارتكب مشكلة في الخارج، أخذته الشرطة، يحتاج ضماناً". تذكر:\textbf{أي مكالمة تطلب تحويلاً مالياً لـ "إطلاق سراح" احتيال}. الشرطة والمحاكم لا تطلب المال هاتفياً.
  \item \textbf{(2) "بطاقة التأمين الطبي خاصتك موقوفة"}— رسالة تقول البطاقة غير طبيعية وتطلب النقر على رابط "للتفعيل". تذكر:\textbf{مكتب التأمين الصحي والضمان الاجتماعي لا يرسل رسائل تطلب النقر على روابط}اتصل بـ 12333 (الخط الرسمي للضمان الاجتماعي) للتحقق.
  \item \textbf{(3) "أرباح استثمارية عالية، مكسب مضمون"}— مجموعات WeChat ورسائل توصي بـ "استثمار تقاعدي" "مشروع دعم وطني"، تضمن عائداً مرتفعاً. تذكر:\textbf{ما يضمن عائداً سنوياً يزيد عن 5\% و"رأس مال مضمون"، 99.9\% احتيال}.
  \item \textbf{(4) "أنا مديرك/صديقك القديم، غيرت رقمي"}— انتحال شخصية معرفة للاقتراض. تذكر:\textbf{الاتصال للتحقق أنفع من أي شيء}من يقول "لا أستطيع الرد على المكالمة" غالباً محتال.
  \item \textbf{(5) "بيض مجاني/محاضرة منتجات صحية"}— أنشطة مجانية تحث كبار السن على حضور "محاضرة صحية" ثم بيع منتجات بسعر مرتفع. تذكر:\textbf{لا شيء مجاني، الأنشطة المجانية وراءها تجارة}.
  \item \textbf{(6) "هاتفك مصاب بفيروس، يحتاج إصلاحاً عن بعد"}— يطلب المحتال تنزيل تطبيق تحكم عن بعد (مثل Sunflower وTeamViewer)، ثم التحكم في هاتفك عن بعد لسرقة المال. تذكر:\textbf{لا تسمح لأحد أبداً بالتحكم في هاتفك عن بعد!}شركات الهواتف ومشغلو الاتصالات والبنوك لا تتصل بك لتثبيت برنامج تحكم عن بعد.
\end{itemize}

\booksubsection{ثلاثة خطوط دفاع، بعد إقامتها تصد 90\% من الاحتيال:}

\begin{enumerate}
  \item \textbf{الخط الأول: اعتراض الهاتف نفسه}— هواتف Android تفتح "اعتراض المكالمات المزعجة" (في تطبيق الاتصال أو تطبيق Mobile Manager)، وiPhone يفتح "كتم المكالمات غير المعروفة" (الإعدادات > الهاتف > كتم المكالمات غير المعروفة).
  \item \textbf{الخط الثاني: اعتراض المشغل}— أرسل رسالة "KTFSR" إلى 10086، لتفعيل اعتراض المكالمات المزعجة عالية التردد من China Mobile مجاناً؛ China Telecom وChina Unicom لديهما خدمات مماثلة. اتصل بخدمة عملاء المشغل لفتحها.
  \item \textbf{الخط الثالث: تطبيق مركز مكافحة الاحتيال الوطني}— نزّل تطبيق "مركز مكافحة الاحتيال الوطني" من متجر التطبيقات، سجّل وفعّل تنبيه المكالمات الواردة. هذا التطبيق مرتبط بقاعدة بيانات الأمن العام، بمجرد اتصال رقم محتال يُحدد.
\end{enumerate}

\begin{tcbwarn}{"نشيد مكافحة الاحتيال" لكبار السن (اطبعه والصقه بجانب الهاتف)}
\begin{itemize}
  \item \textbf{مكالمة غريبة تطلب تحويلاً أغلق!}
  \item \textbf{رسالة برابط تطلب النقر احذف!}
  \item \textbf{من يقول أرباحاً عالية مضمونة احظر!}
  \item \textbf{من يطلب تثبيت برنامج تحكم عن بعد أغلق فوراً!}
  \item \textbf{الشرطة الحقيقية لا تجري قضايا هاتفياً ولا تطلب المال!}
  \item \textbf{إذا لم تكن متأكداً، اتصل بأبنائك أولاً واسأل.}
\end{itemize}

\end{tcbwarn}

\booksection{أربعة: المساعدة عن بعد — في الخارج تستطيع إصلاح هاتف الوالدين}

أحياناً تواجه كبار السن مشكلة، تصفها في الهاتف لفترة طويلة ولا يجدون الزر. هنا تظهر فائدة “المساعدة عن بعد” — يمكنك رؤية شاشة هاتفهم على هاتفك ومساعدتهم. كل علامة لها حل مختلف:

\begin{itemize}
  \item \textbf{Huawei/Honor}: تطبيق "Changlian" مدمج، يحتوي وظيفة "مشاركة الشاشة". ثبّت Changlian على هاتفك وهاتف والديك، أثناء مكالمة فيديو اضغط "مشاركة الشاشة"، يمكنك رسم دوائر على شاشتهم وإخبار والدتك "اضغطي هنا".
  \item \textbf{Xiaomi/Redmi}: تطبيق Xiaomi Video Call يحتوي "مساعدة عن بعد"، يدعم مشاركة الشاشة + مكالمة صوتية متزامنة.
  \item \textbf{OPPO/vivo}: "Family Space" أو "Remote Guard" مدمج، ابحث عن "بعد" في الإعدادات.
  \item \textbf{حل عام (لكل العلامات)}: نزّل تطبيق "Sunflower Remote Control" (النسخة المجانية كافية)، ثبّته على هاتفك وهاتف والديك. سجل بنفس الحساب أو أضف صديقاً، دع والديك يقرؤون رمز التحقق المعروض على هاتفهم، ويمكنك رؤية والتحكم بشاشتهم عن بعد.
  \item \textbf{مستخدمو iPhone}: مكالمات FaceTime تدعم "مشاركة الشاشة". أثناء مكالمة FaceTime، اضغط زر مشاركة الشاشة، ويمكن للطرف الآخر رؤية شاشتك. لكن قيود أمان iOS لا تسمح بالتحكم عن بعد، فقط "تدريس توجيهي".
\end{itemize}

\booksection{خمسة: “تعديل WeChat لكبار السن” — قل هذه الجمل للوالدين}

WeChat أكثر تطبيق يستخدمه كبار السن، لكن واجهته مزدحمة ومعقدة. التالي تعديلات فورية:

\begin{itemize}
  \item \textbf{تكبير خط WeChat نفسه}: WeChat أنا الإعدادات عام حجم الخط، اسحب للأكبر. هذا يؤثر على خط WeChat الداخلي فقط، وليس خط نظام الهاتف.
  \item \textbf{تفعيل "وضع الرعاية"}: WeChat أنا الإعدادات وضع الرعاية تفعيل. بعد التفعيل، تكبر الخطوط والأزرار، ويزداد تباين الألوان، ويمكن بضغطة واحدة "الاستماع للرسائل النصية" (اضغط على رسالة نصية لقراءتها صوتياً).
  \item \textbf{مساعدة الوالدين في ملاحظات جهات الاتصال المهمة}: غيّر أسماء الأبناء والأقارب في دليل العناوين لسهولة التعرف — مثلاً "الابن الأكبر-Xiaoming" "الابنة الثانية-Xiaohong"، ليعرفوا من هو عند فتح WeChat.
  \item \textbf{علّمهم "الإدخال الصوتي" بدلاً من الكتابة}: على لوحة المفاتيح بجانب صندوق الدردشة، اضغط مطولاً على مفتاح المسافة (Android) أو أيقونة الميكروفون (iPhone)، وتحدث مباشرة ليتحول لنص ويُرسل. دقة التعرّف عالية جداً الآن، وهي راحة "تغيير نوعي" لكبار السن الذين لا يعرفون بينين.
  \item \textbf{إغلاق الوظائف غير المستخدمة في صفحة "اكتشاف"}: WeChat أنا الإعدادات عام إدارة صفحة الاكتشاف، أغلق "لحظات" "فيديو" "بث مباشر" "تسوق" "ألعاب" غير المستخدمة، لتقليل التشتت.
\end{itemize}

\booksection{ستة: ماذا لو ضاع الهاتف؟ اصنع “خطة طوارئ” لكبار السن}

الاستعداد يقي من الفوضى. قبل أن يضيع الهاتف، أكمل هذه الأمور:

\begin{enumerate}
  \item \textbf{تأكد من تفعيل "Find My Phone"}: iPhone في الإعدادات اسمك Find Find My iPhone؛ Android ابحث عن "Find Device" في الإعدادات. تأكد أن المفتاح مضيء.
  \item \textbf{دوّن كلمات المرور والحسابات المهمة على ورقة}: كلمات مرور Apple ID / Huawei / Xiaomi، احفظ نسخة في درج مغلق بالمنزل، ونسخة لابن موثوق. بعد فقدان الهاتف تحتاج تسجيل الدخول لهذه الحسابات للقفل والتحديد عن بُعد، نسيان كلمة المرور مشكلة كبيرة.
  \item \textbf{علّم كبار السن جملة واحدة}:\textbf{"إذا ضاع الهاتف، أول شيء اتصل بالأسرة."}يمكن للأسرة تسجيل الدخول فوراً لـ iCloud أو صفحة Find Device في Android، ووضع علامة "وضع الفقدان" عن بُعد، هكذا من يجد الهاتف لا يستطيع فتحه أو استخدامه.
\end{enumerate}

\booksection{سبعة: توافق السماعات اللاسلكية وسماعات الأذن — وظيفة جيدة مهملة}

كثير من كبار السن لديهم مشاكل سمع، لكنهم لا يعرفون أن الهاتف يمكنه العمل مع سماعات الأذن. iPhone وبعض هواتف Android تدعم سماعات الأذن MFi (Made for iPhone)، يمكن نقل صوت المكالمة مباشرة لسماعة الأذن، أوضح بكثير من رفع الهاتف.

\begin{itemize}
  \item \textbf{طريقة الاتصال}: iPhone الإعدادات الوصول المساعد أجهزة السمع، افتح Bluetooth، واتبع تعليمات اقتران سماعة الأذن.
  \item \textbf{إذا لم تستخدم سماعة أذن}: في "الوصول المساعد الصوت/البصري" افتح "تقليل ضجيج المكالمة" و"الصوت الأحادي"، يحسّن وضوح المكالمة.
  \item \textbf{أندرويد}: في "الوصول المساعد تعزيز السمع" إعدادات مماثلة، الموقع يختلف قليلاً حسب العلامة، ابحث عن "سماع".
\end{itemize}

\booksection{ثمانية: الوظائف الصحية — الهاتف ليس مجرد أداة اتصال}

\begin{itemize}
  \item \textbf{SOS طارئ}: اضغط زر الطاقة خمس مرات متتالية (معظم الهواتف)، يتصل الهاتف تلقائياً برقم الطوارئ ويرسل الموقع لجهة الطوارئ المُعدة. اضبطها لكبار السن مسبقاً — iPhone: "الإعدادات SOS طارئ"، Android: "الإعدادات الأمان طلب طوارئ". درّب كبار السن قبل الاستخدام، لتجنب التوتر وعدم الضغط الصحيح.
  \item \textbf{كشف السقوط}: Apple Watch لديه هذه الوظيفة، يكشف السقوط وعدم الحركة لفترة، ويتصل تلقائياً بخدمات الطوارئ ويخطر الأسرة. إذا كانت الميزانية متاحة، وفّر واحدة لكبار السن، قد تنقذ حياة.
  \item \textbf{تذكير تناول الدواء}: تطبيق "الصحة" في iPhone يحتوي وظيفة تذكير الدواء، اضبط اسم الدواء والوقت والجرعة، يظهر تذكير في الوقت المحدد. Android يمكن استخدام "منبه" أو تنزيل تطبيق "تذكير الدواء".
\end{itemize}

\begin{tcbsummary}{خلاصة الفصل}
\begin{itemize}
  \item افتح أولاً الوضع البسيط أو كبّر الخط/الأيقونات يدوياً، ضبط البيئة يحل معظم المشاكل
  \item ثلاث قواعد لمساعدة كبار السن: لا تتعجل، استخدم التشبيهات، اكتب
  \item الست احتيالات الشائعة يجب التذكير بها باستمرار، ثلاثة خطوط دفاع تصد 90\% من الاحتيال
  \item المساعدة عن بعد تتيح لك إصلاح هاتف والديك من الخارج
  \item تفعيل "وضع الرعاية" + "الإدخال الصوتي" في WeChat هما الوظيفتان الأكثر عملية لكبار السن
  \item SOS طارئ وكشف السقوط، إعدادات تنقذ الحياة في اللحظات الحاسمة تستحق خمس دقائق
\end{itemize}

\end{tcbsummary}
